\subsection{Berechnung der Widerstände einer Transistorschaltung}
% Alt: Aufgabe 13
\Aufgabe{
  \label{Aufg13}
  Gegeben ist die abgebildete Transistorschaltung mit den folgenden Parametern: $U_0=\mathrm{5\,V}$ ,
  $U_\mathrm{R1}=\mathrm{0,6\,V}$, $U_\mathrm{L}=\mathrm{15\,V}$, $R_\mathrm{1}=\mathrm{300\,\Omega}$, $R_\mathrm{2}=\mathrm{2\,k\Omega}$ und $R_\mathrm{V}=\mathrm{500\,\Omega}$.
  Damit der gezeichnete Siliziumtransistor schaltet, braucht er einen Basisstrom von $50 \, \mathrm{\mu A}$ um
  durchzuschalten.

  \begin{figure}[H]
    \centering
    \input{Tikz/src/FigTransistorWiderstaende.tex}\alttext{Die Abbildung zeigt eine Transistorschaltung. Links befindet sich eine Spannungsquelle, die eine Spannung \(U_0\) liefert. Von dieser Quelle aus fließt der Strom durch zwei Widerstände, die hintereinander geschaltet sind. Der obere Widerstand ist mit \(R_1\) beschriftet, der untere mit \(R_2\). Die Spannung über dem oberen Widerstand wird als \(U_{\mathrm{R1}}\) angegeben.
      Vom Verbindungspunkt zwischen den beiden Widerständen führt ein weiterer Widerstand \(R_\mathrm{V}\) zu der Basis eines Transistors (ein Halbleiterbauelement mit drei Anschlüssen: Basis, Kollektor und Emitter). An diesem Punkt fließt der Basisstrom \(I_\mathrm{B}\), dargestellt in rot  als Pfeil in Richtung des Transistors.
      Der Transistor hat seinen Emitteranschluss nach unten zum Massepunkt verbunden, durch den der Emitterstrom \(I_\mathrm{E}\) fließt, ebenfalls als Pfeil nach unten dargestellt.
      Auf der rechten Seite des Transistors ist ein Lastwiderstand \(R_\mathrm{L}\) angeschlossen, hinter dem sich eine weitere Spannungsquelle befindet die mit \(U_\mathrm{L}\) gekennzeichnet ist.
      Zusammengefasst handelt es sich um eine typische Verstärkerschaltung mit einem Spannungsteiler am Eingang zur Steuerung des Basisstroms eines Transistors und einem Lastwiderstand auf der Ausgangsseite. Die Ströme im Schaltkreis werden durch rote Pfeile gekennzeichnet, die Spannungen durch blaue Pfeile.}
    \label{fig:FigTransistorWiderstaende}
  \end{figure}

  \begin{itemize}
    \item[a)] Wie groß muss der Widerstand $R_\mathrm{V}$ gewählt werden, damit der Transistor durchschaltet?
    \item[b)] Angenommen der Spannungsabfall des Transistors $U_\mathrm{CE}$ beträgt $7,5 \, \mathrm{V}$, wie groß ist der Strom
          durch $I_\mathrm{E}$?
  \end{itemize}

}


\Loesung{
  \begin{itemize}
    \item[a)]
          \[
            \frac{5\,\text{V}}{2\,\text{k}\Omega + 0,3\,\text{k}\Omega} = \frac{U_1}{300\,\Omega}
            \Rightarrow U_\text{1} = 300\,\Omega \cdot \frac{5\,\text{V}}{2\,\text{k}\Omega + 0,3\,\text{k}\Omega} = 0,65\,\text{V}
          \]
          \[
            5\,\text{V} = 0,65\, \text{V} + U_\text{RV} + 0,6\,\text{V}
          \]
          \[
            U_\text{RV} = 5\,\text{V} - 0,65\,\text{V} - 0,6\,\text{V} = 3,75\,\text{V}
          \]
          \[
            R_\text{V} = \frac{3,75\,\text{V}}{50\,\mu \text{A}} = 75\,\text{k}\Omega
          \]


    \item[b)]
          \[
            I_\text{C} = \frac{15\,\text{V} - 7,5\,\text{V}}{500\,\Omega} = 15\,\text{mA}
          \]
          \[
            I_\text{E} = I_\text{C} + I_\text{B} = 15\,\text{mA} + 50\,\mu \text{A} = 15,05\,\text{mA}
          \]
          \[
            \frac{I_\text{C}}{I_\text{B}} = \beta = 300
          \]
  \end{itemize}
  Siehe: Abschnitt \ref{sec:Bipolartransistor} und \ref{fig:StroemeUndSpannungenBeimBipolartransistor}
}